\clearpage\pagestyle{otMain}\setlength{\otBlocEspacement}{2.5mm}\fontsize{11}{15}\selectfont\corrigesChapitre{3}{Puissances d’exposant positif}
\begin{blocLivre}{cor-3-1}\begin{corrige}{1}{Déplier une puissance}\demarche{Compter les facteurs.}
$5^3=5\times5\times5=\res{125}$. Il y a trois facteurs 5 ; $5\times3$ donnerait 15 et ne représente pas cette puissance.
\end{corrige}\end{blocLivre}
\begin{blocLivre}{cor-3-2}\begin{corrige}{2}{Attention aux parenthèses}\demarche{Repérer la base exacte de la puissance.}
$(-3)^2=(-3)\times(-3)=\res9$. $-3^2=-(3\times3)=\res{-9}$. Enfin $(-3)^3=\res{-27}$, car trois facteurs négatifs donnent un produit négatif.
\end{corrige}\end{blocLivre}
\begin{blocLivre}{cor-3-3}\begin{corrige}{3}{Puissances de dix}\demarche{Compter les zéros donnés par l’exposant.}
$10^5=\res{100\,000}$ : 1 suivi de cinq zéros. $10^1=\res{10}$ et, par convention, $10^0=\res1$.
\end{corrige}\end{blocLivre}
\begin{blocLivre}{cor-3-4}\begin{corrige}{4}{Écrire une puissance}\demarche{Compter les facteurs identiques.}
On compte les facteurs égaux : $\res{3^5}$, $\res{(-2)^3}$ et $\res{7^1}$. Les parenthèses de $(-2)^3$ indiquent que le facteur répété est $-2$.
\end{corrige}\end{blocLivre}
\begin{blocLivre}{cor-3-5}\begin{corrige}{5}{Lire un préfixe}\demarche{Remplacer chaque préfixe par son facteur.}
Kilo signifie $10^3$ : $6$ km $=\res{6\,000\ \mathrm m}$. Giga signifie $10^9$ : 3 Go $=\res{3\times10^9}$ octets, soit trois milliards. Milli signifie $\frac1{10^3}$ : 4 mg $=\frac4{1\,000}$ g $=\res{0{,}004\ \mathrm g}$.
\end{corrige}\end{blocLivre}
\begin{blocLivre}{cor-3-6}\begin{corrige}{6}{Même nombre}\demarche{Utiliser $a^m	imes a^n=a^{m+n}$.}
On additionne les nombres de facteurs : $\res{2^8}$ ; $7=7^1$ donne $\res{7^5}$ ; $6+2+1=9$ donne $\res{10^9}$.
\end{corrige}\end{blocLivre}
\begin{blocLivre}{cor-3-7}\begin{corrige}{7}{Même exposant}\demarche{Utiliser $a^n	imes b^n=(ab)^n$.}
On regroupe les facteurs par paires : $\res{10^4}$ ; $(3\times(-2))^3=\res{(-6)^3}$ ; $(25\times4)^2=100^2=\res{10^4}$.
\end{corrige}\end{blocLivre}
\begin{blocLivre}{cor-3-8}\begin{corrige}{8}{Justifier une règle}\demarche{Compter les facteurs ; calculer chaque membre pour une somme.}
$4^3\times4^2=(4\times4\times4)\times(4\times4)$ contient cinq facteurs 4 : c’est $4^5$. En revanche $4^3+4^2=64+16=\res{80}$, alors que $4^5=1\,024$ : l’égalité est $\res{\text{fausse}}$. La règle concerne les produits.
\end{corrige}\end{blocLivre}
\begin{blocLivre}{cor-3-9}\begin{corrige}{9}{Calculer astucieusement}\demarche{Regrouper des facteurs dont le produit est simple.}
$2^5\times5^5=(2\times5)^5=10^5=\res{100\,000}$. De même $8^3\times125^3=(8\times125)^3=1\,000^3=\res{10^9}$, soit un milliard.
\end{corrige}\end{blocLivre}
\begin{blocLivre}{cor-3-10}\begin{corrige}{10}{Une erreur à corriger}\demarche{Vérifier qu’une propriété s’applique avant de l’utiliser.}
Avec le même nombre 3, on additionne les exposants : $3^2\times3^4=3^6=\res{729}$, alors que $9^6=531\,441$. Pour $2^3\times5^2$, ni les bases ni les exposants ne sont égaux : on calcule $8\times25=\res{200}$.
\end{corrige}\end{blocLivre}
\begin{blocLivre}{cor-3-11}\begin{corrige}{11}{Une quantité qui double}\demarche{Distinguer valeur initiale et facteur répété.}
Chaque étape ajoute un facteur 2. Après six étapes : $3\times2^6=3\times64=\res{192}$. Le facteur initial 3 n’est pas la base de la puissance.
\end{corrige}\end{blocLivre}
\begin{blocLivre}{cor-3-12}\begin{corrige}{12}{Un message transmis}\demarche{Associer chaque étape à un facteur 3.}
À l’étape $k$, le nombre de nouveaux destinataires est $3^k$. À l’étape 5 : $3^5=\res{243}$. Au total : $1+3+9+27+81+243=\res{364}$ personnes. On additionne ici des puissances : aucune règle ne les regroupe en une seule.
\end{corrige}\end{blocLivre}
\begin{blocLivre}{cor-3-13}\begin{corrige}{13}{Des fichiers}\demarche{Choisir la même unité avant de multiplier.}
$2\,400\times2{,}5=\res{6\,000}$ Mo. Comme 1 Go = 1 000 Mo, le total vaut $\res6$ Go, soit $6\times10^9$ octets.
\end{corrige}\end{blocLivre}
\begin{blocLivre}{cor-3-14}\begin{corrige}{14}{Comparer de grandes distances}\demarche{Écrire les deux nombres avec la même puissance de 10.}
$35\times10^7=3{,}5\times10^8$. Comme $4>3{,}5$, $\res{A}$ est la plus grande. La somme vaut $4\times10^8+3{,}5\times10^8=\res{7{,}5\times10^8}$ m : on additionne des multiples de la même puissance de 10.
\end{corrige}\end{blocLivre}
\begin{blocLivre}{cor-3-15}\begin{corrige}{15}{Une somme à ne pas confondre}\demarche{Identifier l’opération principale.}
Pour la somme, $8\times10^4=0{,}8\times10^5$, donc $A=\res{4\times10^5}$. Pour le produit, $3{,}2\times8=25{,}6$ et $10^5\times10^4=10^9$ : $B=25{,}6\times10^9=\res{2{,}56\times10^{10}}$. Une somme ne se traite pas comme un produit.
\end{corrige}\end{blocLivre}
